% !TeX program = tectonic
\documentclass[11pt,a4paper]{article}
\usepackage{fontspec}
\usepackage[russian]{babel}
\babelfont{rm}[Language=Default]{Liberation Serif}
\babelfont[Language=Default]{sf}{Carlito}
\babelfont[Language=Default]{tt}{DejaVu Sans Mono}
\usepackage{amsmath,amssymb,amsthm}
\usepackage{geometry}
\geometry{a4paper, top=2.4cm, bottom=2.4cm, left=2.4cm, right=2.4cm}
\usepackage{booktabs,tabularx,enumitem,xcolor,tcolorbox}
\tcbuselibrary{breakable,skins}
\definecolor{linkblue}{HTML}{1F4E79}
\definecolor{statgreen}{HTML}{2F6B3A}
\definecolor{goldacc}{HTML}{8B7E5A}
\usepackage[colorlinks=true,linkcolor=linkblue,urlcolor=linkblue,unicode=true]{hyperref}
\theoremstyle{plain}
\newtheorem{theorem}{Теорема}
\newtheorem{lemma}[theorem]{Лемма}
\newtheorem{proposition}[theorem]{Предложение}
\newtheorem{corollary}[theorem]{Следствие}
\theoremstyle{definition}
\newtheorem{definition}[theorem]{Определение}
\theoremstyle{remark}
\newtheorem{remark}[theorem]{Замечание}
\newtcolorbox{statusbox}[1][]{colback=statgreen!5,colframe=statgreen!75,
  title={\textbf{Сводка доказанного:} #1},fonttitle=\small,boxrule=0.7pt,arc=2pt,breakable}
\newtcolorbox{databox}[1][]{colback=goldacc!6,colframe=goldacc!80,
  title={\textbf{Данные протокола:} #1},fonttitle=\small,boxrule=0.7pt,arc=2pt,breakable}
\newtcolorbox{verifybox}[1][]{colback=linkblue!5,colframe=linkblue!80,
  title={\textbf{Верификация:} #1},fonttitle=\small,boxrule=0.7pt,arc=2pt,breakable}
\widowpenalty=10000
\clubpenalty=10000
\setlength{\parskip}{0.35em}
\newcommand{\CC}{\mathbb{C}}
\newcommand{\RR}{\mathbb{R}}
\newcommand{\ZZ}{\mathbb{Z}}
\newcommand{\QQ}{\mathbb{Q}}
\newcommand{\Ga}{\Gamma}
\newcommand{\Bfun}{\mathrm{B}}
\newcommand{\im}{\operatorname{Im}}
\newcommand{\re}{\operatorname{Re}}
\newcommand{\lcm}{\operatorname{lcm}}
\newcommand{\SNF}{\operatorname{SNF}}
\newcommand{\Jac}{\operatorname{Jac}}
\newcommand{\rank}{\operatorname{rank}}
\newcommand{\Ind}{\operatorname{Index}}
\newcommand{\disc}{\operatorname{disc}}
\begin{document}
\begin{center}
{\LARGE\bfseries Род кривой Ферма}\\[0.4em]
{\large топологический фундамент циклотомических стендов}\\[0.6em]
{\normalsize Исаев Исхак Хамзатович}\\[0.2em]
{\small Программа <<Динамический принцип>> --- лаборатория \texttt{hodge-laboratory}}\\[0.15em]
{\small Отдельная монография теоремы · Монография 1/16}
\end{center}
\vspace{0.4em}
{\color{linkblue}\hrule height 0.8pt}
\vspace{0.8em}

\section{Постановка}

Кривая Ферма $C_N=\{x^N+y^N=z^N\}\subset\mathbb{P}^2$ --- базовый
объект циклотомических стендов программы. Её топология полностью
определяется одним целым числом $N$, и первый шаг конвейера ---
вычисление рода $g$ --- выполняется чисто топологическими
средствами: ни аналитики, ни спектральной теории не требуется. Род
определяет размерность пространства голоморфных форм, размер
период-матрицы и, в конечном счёте, масштаб всей конструкции.

\begin{lemma}[род кривой Ферма]\label{lem:genus}
Пусть $N\ge 4$. Проекция $\pi_x\colon C_N\to\mathbb{P}^1$ есть
морфизм степени $N$, неразветвлённый над $\infty$ и разветвлённый
ровно над $N$ точками $\zeta_N^k$ ($k=0,\dots,N-1$), причём индекс
ветвления в каждой точке $P_k=(\zeta_N^k,0)$ равен $N$. Род кривой
равен
\[
  g \;=\; \frac{(N-1)(N-2)}{2};
  \qquad
  g(15)=91,\quad g(30)=406 .
\]
\end{lemma}

\begin{proof}
Над произвольной точкой $x_0\in\mathbb{P}^1$ уравнение
$y^N=1-x_0^N$ имеет ровно $N$ решений с учётом кратности, поскольку
многочлен $y\mapsto y^N-(1-x_0^N)$ имеет степень $N$; значит,
$\deg\pi_x=N$. В точке $P_k=(\zeta_N^k,0)$ имеем
$1-x^N=-N\zeta_N^{k(N-1)}(x-\zeta_N^k)+\dots$, то есть локально
$y^N=c\,(x-\zeta_N^k)$ с $c\ne0$: параметром служит $y$, а
$x-\zeta_N^k\sim c^{-1}y^N$, поэтому индекс ветвления равен $N$.
Над $\infty$ ветвления нет: при $z=0$ уравнение $X^N+Y^N=0$ имеет
$N$ различных решений с неразветвлённой проекцией. Формула
Римана--Гурвица даёт
\[
  2g-2 \;=\; N\,(2\cdot 0-2) \;+\; N\,(N-1)
  \;=\; -2N + N^2 - N \;=\; N^2-3N,
\]
откуда $2g=(N-1)(N-2)$ и $g=(N-1)(N-2)/2$. Подстановка $N=15$ и
$N=30$ даёт $91$ и $406$.
\end{proof}

\begin{remark}
Та же формула Плюккера для гладкой плоской кривой степени $m$
даёт $\tfrac{(m-1)(m-2)}2$ --- согласование двух прочтений
(кривая степени $N$ и кривая Ферма) служит встроенным контролем:
для $N=4$ обе дают $3$ (квартика Клейна --- тоже кривая рода $3$,
но \emph{другая} кривая: степени моделей различны).
\end{remark}

\begin{databox}[числа]
$g(15)=\tfrac{14\cdot13}2=91$; период-матрица $91\times182$;
$g(30)=\tfrac{29\cdot28}2=406$; период-матрица $406\times812$;
для сравнения: $g(4)=3$, $g(7)=15$, $g(22)=231$.
\end{databox}

\begin{statusbox}[итог]
Род кривой Ферма доказан из формулы Римана--Гурвица; все компоненты
(степень проекции, точки ветвления, индексы) выписаны явно.
Следствия: перепись характеров $91=1+6+84$ и
$1+6+9+30+84+276=406$ согласована с родом (проверка V1 протокола,
точная целочисленная арифметика).
\end{statusbox}

\begin{verifybox}[как воспроизвести]
\texttt{python3 laboratory.py} $\to$ <<Стенд N=15/30 $\to$ V1
перепись>>: точная арифметика подтвердит суммы $91$ и $406$ и
таблицы $h_d$; Lean: \texttt{genus\_fermat} в
\texttt{HodgeLaboratory.lean} --- машинное доказательство
$(15-1)(15-2)/2=91$ и $(30-1)(30-2)/2=406$.
\end{verifybox}

\vfill
{\color{linkblue}\hrule height 0.6pt}
\vspace{0.5em}
{\small\color{gray}
Лицензия: индивидуальная исключительная лицензия Исаева Исхака Хамзатовича.
Все права на вычисления, формулы и тексты принадлежат автору.
Верификация: \texttt{python3 laboratory.py} (пункт <<Стенды>>/<<Сертификаты>>),
Lean: \texttt{verification/lean/HodgeLaboratory.lean}.
}
\end{document}
